
\documentclass[12pt]{article}
%%%%%%%%%%%%%%%%%%%%%%%%%%%%%%%%%%%%%%%%%%%%%%%%%%%%%%%%%%%%%%%%%%%%%%%%%%%%%%%%%%%%%%%%%%%%%%%%%%%%%%%%%%%%%%%%%%%%%%%%%%%%%%%%%%%%%%%%%%%%%%%%%%%%%%%%%%%%%%%%%%%%%%%%%%%%%%%%%%%%%%%%%%%%%%%%%%%%%%%%%%%%%%%%%%%%%%%%%%%%%%%%%%%%%%%%%%%%%%%%%%%%%%%%%%%%
\usepackage{amsmath}
\usepackage[compat2]{geometry}

\setcounter{MaxMatrixCols}{10}
%TCIDATA{TCIstyle=LaTeX article (bright).cst}

%TCIDATA{OutputFilter=LATEX.DLL}
%TCIDATA{Version=5.50.0.2953}
%TCIDATA{<META NAME="SaveForMode" CONTENT="1">}
%TCIDATA{BibliographyScheme=Manual}
%TCIDATA{Created=Monday, May 23, 2011 14:12:05}
%TCIDATA{LastRevised=Wednesday, December 22, 2021 12:24:14}
%TCIDATA{<META NAME="GraphicsSave" CONTENT="32">}
%TCIDATA{<META NAME="DocumentShell" CONTENT="Exams and Syllabi\SW\Assignment">}

\setlength{\topmargin}{-1.0in}
\setlength{\textheight}{9.25in}
\setlength{\oddsidemargin}{0.0in}
\setlength{\evensidemargin}{0.0in}
\setlength{\textwidth}{6.5in}
\def\labelenumi{\arabic{enumi}.}
\def\theenumi{\arabic{enumi}}
\def\labelenumii{(\alph{enumii})}
\def\theenumii{\alph{enumii}}
\def\p@enumii{\theenumi.}
\def\labelenumiii{\arabic{enumiii}.}
\def\theenumiii{\arabic{enumiii}}
\def\p@enumiii{(\theenumi)(\theenumii)}
\def\labelenumiv{\arabic{enumiv}.}
\def\theenumiv{\arabic{enumiv}}
\def\p@enumiv{\p@enumiii.\theenumiii}
\pagestyle{plain}
\setcounter{secnumdepth}{0}
\parindent=0pt
\input{tcilatex}
\begin{document}


\begin{center}
\begin{equation*}
\FRAME{itbpF}{1.3214in}{0.9496in}{0in}{}{}{Figure}{\special{language
"Scientific Word";type "GRAPHIC";maintain-aspect-ratio TRUE;display
"USEDEF";valid_file "T";width 1.3214in;height 0.9496in;depth
0in;original-width 1.3552in;original-height 0.9669in;cropleft "0";croptop
"1";cropright "1";cropbottom "0";tempfilename
'R45S6500.bmp';tempfile-properties "XPR";}}
\end{equation*}%
\textbf{Electromagnetismo}

\emph{Gu\'{\i}a N}$%
%TCIMACRO{\U{b0}}%
%BeginExpansion
{{}^\circ}%
%EndExpansion
10:$ \emph{Inducci\'{o}n electromagn\'{e}tica}

\textit{Profesor: Arturo\ Fern\'{a}ndez P\'{e}rez}
\end{center}

\begin{enumerate}
\item Un lazo de alambre de lados $L=1$ [m] y $w=10$ [cm] est\'{a} a una
distancia $h=1$ [cm] de una corriente $I(t)=3+10t$ [A], como muestra la Fig.
1. Determine el flujo de campo magn\'{e}tico a trav\'{e}s del lazo, la fem
inducida en el mismo y la direcci\'{o}n de la corriente en el lazo. \textbf{%
R: }$\phi _{B}(t)=1,439\times 10^{-6}+\left( 4,8\times 10^{-6}\right) t$
[Wb]; $\varepsilon =-4,8$ [$\mu V$]; Corriente en sentido antihorario.%
\begin{equation*}
\FRAME{itbpFU}{2.3661in}{1.4736in}{0in}{\Qcb{Figura 1}}{}{Figure}{\special%
{language "Scientific Word";type "GRAPHIC";maintain-aspect-ratio
TRUE;display "USEDEF";valid_file "T";width 2.3661in;height 1.4736in;depth
0in;original-width 3.8735in;original-height 2.4025in;cropleft "0";croptop
"1";cropright "1";cropbottom "0";tempfilename
'R45S6501.wmf';tempfile-properties "XPR";}}
\end{equation*}

\item Una bobina de 15 vueltas y radio $R=10$ [cm] rodea un solenoide
extenso de radio $r=2$ [cm] y 1000 vueltas/metro (Fig. 2). Si la corriente
en el solenoide es $I(t)=5\sin (120t)$ [A], determine la fem inducida en la
bobina. \textbf{R: }$\varepsilon =-14\cos (120t)$ [mV] 
\begin{equation*}
\FRAME{itbpFU}{1.6648in}{1.2142in}{0in}{\Qcb{Figura 2}}{}{Figure}{\special%
{language "Scientific Word";type "GRAPHIC";maintain-aspect-ratio
TRUE;display "USEDEF";valid_file "T";width 1.6648in;height 1.2142in;depth
0in;original-width 3.7619in;original-height 2.7363in;cropleft "0";croptop
"1";cropright "1";cropbottom "0";tempfilename
'R45S6502.wmf';tempfile-properties "XPR";}}
\end{equation*}

\item Determine la corriente a trav\'{e}s de la secci\'{o}n $PQ$ de longitud 
$a=65$ [cm] de la Fig. 3. El circuito es colocado en una regi\'{o}n de campo
magn\'{e}tico cuya magnitud es $B(t)=0.001t.$ La resistencia por unidad de
longitud del cable es $0.1$ [$\Omega /m$]. \textbf{R:} $I_{PQ}=0,85$ [mA]%
\begin{equation*}
\FRAME{itbpFU}{2.322in}{0.9937in}{0in}{\Qcb{Figura 3}}{}{Figure}{\special%
{language "Scientific Word";type "GRAPHIC";maintain-aspect-ratio
TRUE;display "USEDEF";valid_file "T";width 2.322in;height 0.9937in;depth
0in;original-width 4.8724in;original-height 2.0695in;cropleft "0";croptop
"1";cropright "1";cropbottom "0";tempfilename
'R45S6503.wmf';tempfile-properties "XPR";}}
\end{equation*}

\item Un lazo rectangular de 40 x 30 [cm] se mueve a una rapidez $v=2$
[cm/s] en una regi\'{o}n de campo magn\'{e}tico uniforme $B=1,25$ [T]. La
regi\'{o}n con campo magn\'{e}tico termina abruptamente, como se muestra en
la Fig. 4. Determine la fem inducida cuando el lazo est\'{a}

\begin{enumerate}
\item Totalmente dentro del campo. \textbf{R: }$\varepsilon =0$ [V]

\item Parcialmente dentro del campo. \textbf{R: }$\varepsilon =-10$ [mV]

\item Totalmente fuera del campo. \textbf{R:} $\varepsilon =0$ [V]%
\begin{equation*}
\FRAME{itbpFU}{1.932in}{1.0957in}{0in}{\Qcb{Figura 4}}{}{Figure}{\special%
{language "Scientific Word";type "GRAPHIC";maintain-aspect-ratio
TRUE;display "USEDEF";valid_file "T";width 1.932in;height 1.0957in;depth
0in;original-width 4.6501in;original-height 2.6247in;cropleft "0";croptop
"1";cropright "1";cropbottom "0";tempfilename
'R45S6604.wmf';tempfile-properties "XPR";}}
\end{equation*}
\end{enumerate}

\item Suponga que el lazo de corriente de la Fig. 5 rota alrededor del eje $%
y $; el eje $x$; el eje $z$. \textquestiondown Cu\'{a}l es el valor m\'{a}%
ximo de la fem inducida en cada caso si $A=600$ [cm$^{2}$], $\omega =35$
[rad/s] y $B=0.45$ [T]? \textbf{R: }Alrededor del eje $y$ y del eje $z,$ $%
\varepsilon _{\max }=0,945$ [V]. Alrededor del eje $x$, $\varepsilon _{\max
}=0$ [V].%
\begin{equation*}
\FRAME{itbpFU}{1.6198in}{1.3111in}{0in}{\Qcb{Figura 5}}{}{Figure}{\special%
{language "Scientific Word";type "GRAPHIC";maintain-aspect-ratio
TRUE;display "USEDEF";valid_file "T";width 1.6198in;height 1.3111in;depth
0in;original-width 3.0952in;original-height 2.5002in;cropleft "0";croptop
"1";cropright "1";cropbottom "0";tempfilename
'R45S6605.wmf';tempfile-properties "XPR";}}
\end{equation*}

\item En la Fig. 6 se muestran las secciones transversales de dos solenoides
muy largos y la direcci\'{o}n de sus campos magn\'{e}ticos, que se
incrementan a una tasa de $2$ [$\frac{T}{s}]$. Utilizando las leyes de
Kirchhoff, determine la corriente inducida en cada resistencia. \textbf{R: }$%
I_{6\Omega }=1,286$ [mA], $I_{5\Omega }=17,14$ [mA], $I_{3\Omega }=18,43$
[mA]%
\begin{equation*}
\FRAME{itbpFU}{2.5434in}{1.1986in}{0in}{\Qcb{Figura 6}}{}{Figure}{\special%
{language "Scientific Word";type "GRAPHIC";maintain-aspect-ratio
TRUE;display "USEDEF";valid_file "T";width 2.5434in;height 1.1986in;depth
0in;original-width 5.9828in;original-height 2.8055in;cropleft "0";croptop
"1";cropright "1";cropbottom "0";tempfilename
'R45S6606.wmf';tempfile-properties "XPR";}}
\end{equation*}

\item Un disco de aluminio de radio $r_{1}=5\,\ [cm]$ y resistencia $%
R=0.0003 $ [$\Omega $] se ubica sobre un solenoide de 1000 vueltas/metro y
radio $r_{2}=3$ [cm], como muestra la Fig. 7. Asuma que en el campo del
solenoide en el extremo superior es la mitad que en el centro del mismo. Si
la corriente en el solenoide aumenta a una tasa de $25$ $[\frac{A}{s}]$ 
\textquestiondown Cu\'{a}l es la corriente inducida en el anillo? \textbf{R: 
}$I=0,148$ [A] \textquestiondown Cu\'{a}l es el campo magn\'{e}tico
producido por la corriente inducida, en el centro del anillo? \textbf{R: }$%
B=1,86\times 10^{-6}$ [T]%
\begin{equation*}
\FRAME{itbpFU}{1.1632in}{1.8092in}{0in}{\Qcb{Figura 7}}{}{Figure}{\special%
{language "Scientific Word";type "GRAPHIC";maintain-aspect-ratio
TRUE;display "USEDEF";valid_file "T";width 1.1632in;height 1.8092in;depth
0in;original-width 3.5959in;original-height 5.6109in;cropleft "0";croptop
"1";cropright "1";cropbottom "0";tempfilename
'R45S6607.wmf';tempfile-properties "XPR";}}
\end{equation*}

\item En la Fig. 8 un im\'{a}n se dirige hacia un lazo de corriente 
\textquestiondown La diferencia de potencial $V_{a}-V_{b},$ es positiva,
negativa o cero? \textbf{R: }Negativa, ya que la corriente inducida circula
desde $b$ hacia $a.$%
\begin{equation*}
\FRAME{itbpFU}{1.9545in}{1.4883in}{0in}{\Qcb{Figura 8}}{}{Figure}{\special%
{language "Scientific Word";type "GRAPHIC";maintain-aspect-ratio
TRUE;display "USEDEF";valid_file "T";width 1.9545in;height 1.4883in;depth
0in;original-width 6.3728in;original-height 4.8473in;cropleft "0";croptop
"1";cropright "1";cropbottom "0";tempfilename
'R45S6608.wmf';tempfile-properties "XPR";}}
\end{equation*}

\item El inductor de la Fig. 9 tiene una inductancia $L=260$ [$mH]$ y es
atravesado por una corriente $i$. Si la corriente decae con una tasa $\frac{%
di}{dt}=-0,0180$ [A/s]$,$ encuentre la fem inducida en el inductor y
determine qu\'{e} punto (a \'{o} b$)$ est\'{a} a un potencial m\'{a}s alto. 
\textbf{R: }$\varepsilon =4,68$ [mV]$.$ El punto $b$ est\'{a} a mayor
potencial, ya que la corriente $i$ disminuye en el tiempo y la corriente
inducida va en la misma direcci\'{o}n.%
\begin{equation*}
\FRAME{itbpFU}{2.2373in}{1.0482in}{0in}{\Qcb{Figura 9}}{}{Figure}{\special%
{language "Scientific Word";type "GRAPHIC";maintain-aspect-ratio
TRUE;display "USEDEF";valid_file "T";width 2.2373in;height 1.0482in;depth
0in;original-width 3.4134in;original-height 1.5852in;cropleft "0";croptop
"1";cropright "1";cropbottom "0";tempfilename
'R45S7N0C.wmf';tempfile-properties "XPR";}}
\end{equation*}

\item En la Fig. 10, se muestra el gr\'{a}fico de la corriente $i(t)$ a trav%
\'{e}s de un inductor de $4,6$ [H] que tiene una resistencia $R=12$ [$\Omega 
$]. Si la escala del gr\'{a}fico llega hasta 8,0 [A] y t$_{s}=6$ [ms],
determine la magnitud de la corriente inducida a trav\'{e}s del inductor, en
los intervalos (a) 0 a 2 ms (b) 2 a 5 ms y (c) 5 a 6 ms.%
\begin{equation*}
\FRAME{itbpFU}{2.9248in}{2.0487in}{0in}{\Qcb{Figura 10}}{}{Figure}{\special%
{language "Scientific Word";type "GRAPHIC";maintain-aspect-ratio
TRUE;display "USEDEF";valid_file "T";width 2.9248in;height 2.0487in;depth
0in;original-width 7.504in;original-height 5.2442in;cropleft "0";croptop
"1";cropright "1";cropbottom "0";tempfilename
'R45SAH0F.wmf';tempfile-properties "XPR";}}
\end{equation*}

\item Un lazo rectangular de $N=50$ vueltas ($b=5$ [cm] y $l=15$ [cm]) se
encuentra a una distancia $a=3$ [cm] de un conductor recto muy largo, por cu%
\'{a}l circula una corriente $i$ (Fig. 11). Determine la inductancia mutua
de este sistema. \textbf{R: }$M=1,47$ [$\mu H$]%
\begin{equation*}
\FRAME{itbpFU}{2.7432in}{2.0297in}{0in}{\Qcb{Figura 11}}{}{Figure}{\special%
{language "Scientific Word";type "GRAPHIC";maintain-aspect-ratio
TRUE;display "USEDEF";valid_file "T";width 2.7432in;height 2.0297in;depth
0in;original-width 5.2088in;original-height 3.8475in;cropleft "0";croptop
"1";cropright "1";cropbottom "0";tempfilename
'R45S8S0E.wmf';tempfile-properties "XPR";}}
\end{equation*}
\end{enumerate}

\end{document}
